% Copyright 2026 The terCAD team. All rights reserved.
% Use of this content is governed by a CC BY-NC-ND 4.0 license that can be found in the LICENSE file.

\noindent \textbf{Fingrom} is an open source code, cross-platform (iOS, Android, macOS,
Windows, Linux, Web) financial accounting application without ads or limitations. The goal is
for it to be intuitive, efficient and inclusive, enabling anyone to effortlessly manage their
personal finances, regardless of the device they use, the language they speak or their level
of financial expertise.

This manual guides you through using the application, from the initial launch to accessing
its most advanced features. These include configuring accounts and budgets, recording
everyday transactions, setting goals, analysing metrics, synchronising devices and
protecting data.


\subsection*{Why Fingrom}

Most finance trackers either restrict access to key features unless users pay for a
subscription, display advertisements to monetise users' attention, or require a cloud account
that stores sensitive financial history on someone else's servers.

\textbf{Fingrom} operates according to a different set of principles: every feature is
available to all users free of charge; sponsorship only supports the further development
of the application and never unlocks additional features; there are no advertisements, so
the entire screen is dedicated to your finances; and no account registration is required
because records are stored locally on the device, can be encrypted and synchronized directly
between your own devices (peer-to-peer), and backed up wherever you choose -- to a file or
your own WebDAV storage. In short, you have full control over both the application and your
data.


\subsection*{Core Concepts}

Fingrom is built around a small number of entities. Understanding how these entities relate
to each other will make the rest of the manual easier to follow.

\begin{description}
  \item[Accounts] are places where money is kept, such as cash, bank and debit cards, credit
  cards, deposits, loans, investments and cryptocurrency wallets. Each account has its own
  currency and balance and logs all the transactions that affect it.
  \item[Budget Categories] describe what money is spent on, such as groceries, rent,
  transport and leisure. A category can have a fixed limit, a limit configured per month, a
  limit relative to income, or no limit (in which case only the amount spent is tracked).
  \item[Transactions] change the balance of accounts and budgets.
    \begin{itemize}
      \item \emph{Bills} (expenses) decrease an account and are attributed to a budget
      category. A bill can be split across several categories.
      \item \emph{Incomes} (invoices) increase an account;
      \item \emph{Transfers} move money between two accounts, including accounts in different
      currencies.
    \end{itemize}
  \item[Goals] are the targets you are saving for, with a desired amount and deadline.
  \item[Currencies] and their exchange rates allow you to mix accounts in different
  currencies while still seeing the overall picture in a single default currency.
\end{description}

\noindent You can group accounts and budget categories on the main page simply by adding the
\q{/} symbol to their names. For example, \q{Bank/Debit} and \q{Bank/Savings} are displayed
together under \q{Bank}.


\subsection*{Capabilities at a Glance}

Beyond everyday bookkeeping, the application provides a range of additional features:

\begin{itemize}
  \item \textbf{Automation} -- recurring payments created automatically on schedule (with a
  home screen widget on Android) and budget category prediction for a new bill.
  \item \textbf{Metrics} -- a budget forecast based on a Monte Carlo simulation; budget
  limits versus monthly expenses; a candlestick (OHLC) chart of accounts; an income health
  radar; currency distribution; year-to-date expenses; a category bar chart; a goals gauge;
  and a historical currency chart;
  \item \textbf{Import and export} -- importing and exporting bills and incomes from CSV, QIF
  and OFX files exported by banks or other applications, and exporting data to an Excel XLSX
  workbook.
  \item \textbf{Synchronization and recovery} -- peer-to-peer synchronisation between devices
  and backup or restoration via a WebDAV server or direct file.
  \item \textbf{Security} -- data encryption, protection by a one-time password code, a
  recovery code and biometric authentication.
  \item \textbf{Personalization} -- a configurable main page with separate layouts for
  different screen sizes, a dark, light or system theme with a custom colour palette, the
  start day of the week and month, and the budget timeline (weekly, monthly or yearly).
\end{itemize}


\subsection*{Supported Platforms}

Fingrom is distributed through the official stores of each platform, which can be found on
the relevant app stores.

\begin{center}
\begin{tabularx}{\textwidth}{|l|X|}
  \hline
  \textbf{Platform} & \textbf{Distribution} \\
  \hline
  iOS, macOS & Apple App Store \\
  \hline
  Android & Google Play, Galaxy Store, Huawei AppGallery, Amazon Appstore \\
  \hline
  Linux & Snap Store, Flathub, AppImage \\
  \hline
  Windows & Microsoft Store \\
  \hline
  Web & Any modern browser, no installation required \\
  \hline
\end{tabularx}
\end{center}

\noindent Every release is also published on GitHub for users who prefer to install the
application manually.

\begin{itemize}[leftmargin=*]

\item \q{Fingrom} -- open source code financial accounting app\\
\href{https://github.com/lyskouski/app-finance}{\small https://github.com/lyskouski/app-finance}

\item Apple Store (iOS, macOS):\\
\href{https://apps.apple.com/us/app/fingrom/id6463955600}{\small https://apps.apple.com/us/app/fingrom/id6463955600}

\item Google Play (Android):\\
\href{https://play.google.com/store/apps/details?id=com.tercad.fingrom}{\small https://play.google.com/store/apps/details?id=com.tercad.fingrom}

\item Galaxy Store (Android):\\
\href{https://galaxystore.samsung.com/detail/com.tercad.fingrom}{\small https://galaxystore.samsung.com/detail/com.tercad.fingrom}

\item Huawei App Gallery (Android):\\
\href{https://appgallery.huawei.com/\#/app/C109437079}{\small https://appgallery.huawei.com/\#/app/C109437079}

\item Amazon Appstore (Android):\\
\href{https://www.amazon.com/gp/product/B0DKFNJ45R}{\small https://www.amazon.com/gp/product/B0DKFNJ45R}

\item Linux (Snap Store):\\
\href{https://snapcraft.io/fingrom}{\small https://snapcraft.io/fingrom}

\item Linux (Flathub):\\
\href{https://flathub.org/apps/com.tercad.fingrom}{\small https://flathub.org/apps/com.tercad.fingrom}

\item Linux (AppImage):\\
\href{https://www.appimagehub.com/p/2173504}{\small https://www.appimagehub.com/p/2173504}

\item Windows:\\
\href{https://apps.microsoft.com/detail/fingrom/9NNPDJ2ST0HV}{\small https://apps.microsoft.com/detail/fingrom/9NNPDJ2ST0HV}

\item Web:\\
\href{https://lyskouski.github.io/app-finance/}{\small https://lyskouski.github.io/app-finance/}

\end{itemize}


\subsection*{Who This Book Is For}

This manual is intended for anyone who uses or is considering using Fingrom:

\begin{itemize}
  \item \textbf{newcomers} who want to set up the application properly from the outset and
  develop the habit of tracking their finances.
  \item \textbf{users migrating} from spreadsheets or other applications who need to import
  their history and preserve existing balances.
  \item \textbf{experienced users} who want to make the most of budgeting strategies,
  automation, metrics and synchronization across several devices.
\end{itemize}

\noindent No prior knowledge of accounting is required. Financial terms are explained as they
first appear.


\subsection*{Feedback and Contribution}

Fingrom is developed openly, which is a process that is open to anyone who wishes to
participate. If you find anything in the application or this book that is unclear, incorrect
or missing, please open an issue at
\href{https://github.com/lyskouski/app-finance/issues}{https://github.com/lyskouski/app-finance/issues}.
Translations, bug reports, ideas and pull requests are always welcome and are gratefully
acknowledged in the application's \textbf{About} section under the \textbf{Contributors}
heading.
